\section{Computational Experiments}
\label{sec:experiments}

The computational study uses component-wise comparisons to isolate each of the
three design refinements introduced in Section~\ref{sec:audit}.
In each comparison, exactly one component differs between the two variants,
so the contribution of that specific modification is measured without confounding.
This structure directly mirrors the analysis in Section~\ref{sec:audit}: each
comparison targets one of the three design aspects in isolation.

\subsection{Experimental Setup}

\paragraph{Instances.}
We use the large-scale MSCFLP benchmark families of
\citet{avella2008effective}, also used by \citet{guastaroba2012kernel}.
The main study covers the public Test Bed~A and Test Bed~C families listed in
Table~\ref{tab:instance_families}.
Test Bed~A and Test Bed~C share the same five facility-client scales but use
independently drawn instances, providing a robustness check across instance
variation within each scale.

\begin{table}[htbp]
\centering
\caption{Large-scale benchmark families used in the computational study.}
\label{tab:instance_families}
\small
\setlength{\tabcolsep}{5pt}
\begin{tabular}{llrrrr}
\toprule
Test Bed & Family & $|\J|$ & $|\I|$ & $N$ & Instance indices \\
\midrule
A & p800-4400  &   800 & 4{,}400 & 30 & 1--30 \\
A & p1000-1000 & 1{,}000 & 1{,}000 & 30 & 1--30 \\
A & p1000-4000 & 1{,}000 & 4{,}000 & 30 & 1--30 \\
A & p1200-3000 & 1{,}200 & 3{,}000 & 30 & 1--30 \\
A & p2000-2000 & 2{,}000 & 2{,}000 & 30 & 1--30 \\
\midrule
C & p1000-1000 & 1{,}000 & 1{,}000 & 30 & 61--90 \\
\bottomrule
\end{tabular}
\end{table}

\paragraph{Algorithms compared.}
We report four variants, all implemented in Python~3.11 with Gurobi~13.0.2,
run on an Apple M3~Pro (12-core CPU, 18~GB RAM).
Each instance is solved in an independent process; only one instance runs at a
time, so the full memory and CPU are available to each run.
\begin{itemize}
  \item \textbf{KS-LP}: our re-implementation of G\&S~2012 using the full LP
        relaxation (Gurobi simplex).
        All parameters follow G\&S~2012 defaults.
        KS-LP is the within-study baseline; comparing against it measures
        the combined effect of all three fixes.
  \item \textbf{KS-CG}: the same KS-LP restricted-MILP phase with the full LP
        replaced by column generation.
        Comparing KS-CG against KS-LP isolates Fix~1 (LP cost) alone.
  \item \textbf{KS-CG-WS}: KS-CG with feasibility-preserving stage transitions
        and warm starting added, but retaining the original $\gamma$-threshold
        arc selection.
        Comparing KS-CG-WS against KS-CG isolates Fix~3 alone.
  \item \textbf{FKS-CG-WS}: Funnel Kernel Search --- KS-CG-WS with the
        $\gamma$-threshold replaced by the flat-$k$ criterion,
        $(k_1, k_2, k_3) = (50, 20, 10)$.
        Comparing FKS-CG-WS against KS-CG-WS isolates Fix~2 alone.
        Comparing FKS-CG-WS against KS-LP gives the combined end-to-end picture.
\end{itemize}

\paragraph{Implementation details.}
KS-LP uses the parameter choices of G\&S~2012: the LP-defined initial kernel,
bucket size $l_{\text{buck}} = |\K|$, at most three analyzed buckets ($\bar{N}_B = 3$),
objective cutoffs, and the bucket-forcing constraint.
FKS uses the same kernel and bucket logic but replaces the single restricted-MILP
phase with the flat-$k$ funnel described in Section~\ref{sec:algorithm}.
The restricted-MILP time limit is 300~seconds per subproblem for all variants.
G\&S~2012 allocated 900~seconds per subproblem~\citep{guastaroba2012kernel}
(Section~3.2), noting that ``the software often finds quickly an optimal
or near-optimal solution'' and that the limit exists mainly as a safeguard.
Modern hardware and Gurobi reduce per-subproblem solve times considerably
relative to the CPLEX~12.2 environment of G\&S~2012, so 300~seconds proves
sufficient: our KS-LP recreation matches or slightly exceeds their published
aggregate results on p1000-1000 (Section~\ref{sec:literature}).
Because all four variants share this 300-second limit, all within-study
comparisons are unaffected by the difference from the original paper.
Runs are recorded with instance name, test bed, LP solver, stage widths,
LP time, restricted-MILP time, objective value, and gap.

\paragraph{Metrics.}
The LP gap is
\[
  \text{gap}(\%) = \frac{z^H - z^{\mathrm{LP}}}{z^{\mathrm{LP}}} \times 100,
\]
where $z^H$ is the heuristic objective and $z^{\mathrm{LP}}$ is the LP lower
bound from the same run.
Because our LP lower bounds can differ slightly from those reported by
G\&S~2012 due to solver and hardware differences (see Section~\ref{sec:lp_bounds}),
we also report integer objectives for direct literature comparisons.
Win/tie/loss counts use a tolerance of 0.01 cost units.
Runtime is decomposed into LP time, restricted-MILP time, and total time
whenever the breakdown is relevant.

\subsection{Study 1: LP Bottleneck Removal}
\label{sec:study_lp}

The first study measures how much of the KS setup time is consumed by the LP
phase and how accurately column generation replicates the full LP solution.
This is primarily a timing and fidelity check, not a heuristic-quality
comparison.

Table~\ref{tab:lp_bottleneck} summarizes LP timing across all benchmark families.
Two patterns stand out.
First, CG LP always matches the full LP objective to machine precision: the
two lower bounds agree in all reported significant figures on every completed
instance, confirming that column generation supplies the same information to KS
without altering its guidance.
Second, the CG speedup grows sharply with instance size.
For p1000-1000 (Test Bed~A), the full LP averages 58.0~seconds and CG reduces
this to 1.3~seconds (45.6$\times$) using 3.4~iterations and activating only 3.3\%
of the $10^6$ possible shipping arcs.
For p2000-2000, the full LP averages 710.7~seconds and CG reduces this to
4.2~seconds, a 170.2$\times$ speedup.
The active-arc percentage falls as the instance grows (3.8\% for p800-4400 to
2.4\% for p2000-2000), consistent with larger instances having sparser LP
solutions and therefore more redundant columns to avoid.

\begin{table}[htbp]
\centering
\caption{LP relaxation: full LP (Gurobi simplex) versus column-generation LP.
         Speedup is the ratio of average full-LP time to average CG-LP time.
         Active arcs reports the percentage of all $|\J|\times|\I|$ arcs
         variables present in the terminal CG LP.
         CG LP objectives match full LP objectives to machine precision on all
         completed instances.}
\label{tab:lp_bottleneck}
\footnotesize
\setlength{\tabcolsep}{3pt}
\begin{tabular}{llrrrrr}
\toprule
Test Bed & Family & Full LP (s) & CG LP (s) & Speedup & CG iters & Active arcs (\%) \\
\midrule
A & p800-4400  & 276.5 & 5.9 & 46.7$\times$  & 2.6 & 3.8 \\
A & p1000-1000 &  58.0 & 1.3 & 45.6$\times$  & 3.4 & 3.3 \\
A & p1000-4000 & 409.1 & 6.9 & 59.3$\times$  & 2.8 & 3.0 \\
A & p1200-3000 & 389.3 & 4.5 & 85.8$\times$  & 3.3 & 2.5 \\
A & p2000-2000 & 710.7 & 4.2 & 170.2$\times$ & 3.9 & 2.4 \\
\midrule
C & p1000-1000 &   8.8 & 0.7 & 13.5$\times$  & 1.3 & 3.0 \\
\bottomrule
\end{tabular}
\end{table}

\subsubsection{On LP Lower-Bound Differences Between Studies}
\label{sec:lp_bounds}

Our full LP lower bounds for p1000-1000 match those of G\&S~2012 exactly on
tight-capacity instances (instances~1--5), but are slightly higher on
looser-capacity instances (instances~16--30), by up to 25~cost units.
We treat these differences as numerical and reporting differences across
software generations rather than as an algorithmic improvement in the LP bound.
The original study used a different solver environment and reports rounded
values, while our within-study comparisons use LP bounds computed in the same
Gurobi environment for all three variants.
The important validation for this paper is therefore internal: the CG LP
reproduces our full LP lower bound to within $1.3 \times 10^{-5}$ on all
instances, confirming that column generation supplies the same LP information
used by the full-LP KS baseline.

\subsection{Study 2: Restricted-MILP Acceleration}
\label{sec:study_milp}

The second study compares KS-CG and FKS-CG-WS on the restricted-MILP phase.
Both variants receive identical CG LP output, so any difference in MILP time
or solution quality is attributable to the flat-$k$ arc selection and
feasibility-preserving stage transitions combined.
Table~\ref{tab:kscg_fkscg} reports results across all families; a separate
component-wise breakdown isolating each fix individually is given in
Section~\ref{sec:study_components}.

FKS-CG-WS never loses on solution quality: across all five Test Bed~A families
(150 instances total), FKS-CG-WS matches or improves KS-CG on every instance.
The quality advantage is most pronounced on p1000-4000 (23W/7T/0L) and
p2000-2000 (22W/8T/0L), where the flat-$k$ arc selection appears to
consistently deliver better facility combinations than the $\gamma$-threshold.

MILP-phase speedup varies by family and depends on how incumbents interact with
stage transitions.
On families where LP solutions are relatively concentrated (p1000-1000: 1.86$\times$;
p2000-2000: 2.18$\times$), the funnel's narrowing edge sets reduce subproblem size
at each stage and FKS-CG-WS is substantially faster.
On families with more diffuse LP support (p800-4400: 0.94$\times$; p1000-4000:
0.80$\times$), incumbent-arc preservation fills later-stage edge sets with arcs
from the Stage~1 solution, limiting the thinning effect of the funnel.
In these cases FKS-CG-WS's MILP phase is slightly slower than KS-CG, but end-to-end
FKS-CG-WS remains faster once the LP-phase savings are included (Table~\ref{tab:end_to_end}).

Test Bed~C p1000-1000 (30 independent instances, indices 61--90) shows
1.76$\times$ MILP speedup with 7W/22T/1L.
The single loss is instance~82, where KS-CG finds an objective 2.6 cost units
lower (a 0.00007\% relative difference), consistent with numerical noise rather
than a systematic quality difference.

\begin{table}[htbp]
\centering
\caption{Restricted-MILP phase: KS-CG versus FKS-CG-WS on the same CG LP.
         MILP time excludes the shared LP phase.
         Speedup is KS-CG time divided by FKS-CG-WS time.
         W/T/L counts wins, ties, and losses for FKS-CG-WS relative to KS-CG
         by integer objective ($\pm 0.01$ tolerance).}
\label{tab:kscg_fkscg}
\footnotesize
\setlength{\tabcolsep}{3.5pt}
\begin{tabular}{llrrrrr}
\toprule
 & & & \multicolumn{2}{c}{Avg.\ MILP time (s)} & & \\
\cmidrule(lr){4-5}
TB & Family & $N$ & KS-CG & FKS-CG-WS & Speedup & W/T/L \\
\midrule
A & p800-4400  & 30 & 513.7 & 546.3 & 0.94$\times$ & 12/18/0 \\
A & p1000-1000 & 30 &  80.6 &  43.4 & 1.86$\times$ &  1/29/0 \\
A & p1000-4000 & 30 & 444.3 & 555.0 & 0.80$\times$ & 23/7/0  \\
A & p1200-3000 & 30 & 508.8 & 481.3 & 1.06$\times$ & 16/14/0 \\
A & p2000-2000 & 30 & 479.2 & 219.5 & 2.18$\times$ & 22/8/0  \\
\midrule
C & p1000-1000 & 30 & 125.1 &  71.2 & 1.76$\times$ &  7/22/1 \\
\bottomrule
\end{tabular}
\end{table}

\subsection{Study 3: End-to-End Performance}
\label{sec:study_end_to_end}

Table~\ref{tab:end_to_end} reports end-to-end results across all families.
FKS-CG-WS never loses on solution quality against KS-LP across all 150 Test Bed~A
instances: the W/T/L counts show wins or ties only.
The same holds against KS-CG except for the single numerically negligible
loss on Test Bed~C instance~82 noted above.

The end-to-end speedup over KS-LP grows with instance size, driven primarily
by the LP phase.
For p1000-1000, where the full LP is already modest (58~s), the combined speedup
is 3.08$\times$.
For p2000-2000, where the full LP averages 710.7~s and column generation
reduces this to 4.2~s, the combined speedup reaches 5.32$\times$.

The pattern across families also illustrates the interaction between the two
phases.
On p800-4400, FKS-CG-WS's MILP phase is slightly slower than KS-CG (0.94$\times$),
but the LP-phase saving of 270~s per instance still produces a 1.52$\times$
end-to-end gain over KS-LP.
On p1000-4000, FKS-CG-WS's MILP phase is again slower than KS-CG (0.80$\times$),
yet end-to-end FKS-CG-WS is 2.01$\times$ faster than KS-LP because the LP saving
(409~s full LP versus 6.9~s CG) dominates.
On p2000-2000, both phases favor FKS-CG-WS simultaneously: the LP phase contributes
a 170$\times$ speedup and the MILP phase contributes a further 2.18$\times$,
yielding the largest combined gain in the study.

\begin{table}[htbp]
\centering
\caption{End-to-end total time (LP + restricted-MILP) for all three variants.
         Speedup is KS-LP time divided by FKS-CG-WS time.
         W/T/L counts wins, ties, and losses for FKS-CG-WS relative to each
         baseline by integer objective ($\pm 0.01$ tolerance).
         Solution quality (gap) is reported separately in
         Table~\ref{tab:kscg_fkscg}.}
\label{tab:end_to_end}
\footnotesize
\setlength{\tabcolsep}{3pt}
\begin{tabular}{llrrrrrrr}
\toprule
 & & & \multicolumn{3}{c}{Avg.\ total time (s)} & & \multicolumn{2}{c}{W/T/L (FKS-CG-WS vs)} \\
\cmidrule(lr){4-6}\cmidrule(lr){8-9}
TB & Family & $N$ & KS-LP & KS-CG & FKS-CG-WS & Speedup & KS-LP & KS-CG \\
\midrule
A & p800-4400  & 30 &   838.3 & 519.6 & 552.3 & 1.52$\times$ & 15/15/0 & 12/18/0 \\
A & p1000-1000 & 30 &   137.9 &  81.8 &  44.7 & 3.08$\times$ &  4/26/0 &  1/29/0 \\
A & p1000-4000 & 30 & 1{,}131.7 & 449.9 & 561.9 & 2.01$\times$ & 17/13/0 & 23/7/0  \\
A & p1200-3000 & 30 & 1{,}048.6 & 513.3 & 485.8 & 2.16$\times$ & 16/14/0 & 16/14/0 \\
A & p2000-2000 & 30 & 1{,}189.3 & 483.6 & 223.7 & 5.32$\times$ & 22/8/0  & 22/8/0  \\
\midrule
C & p1000-1000 & 30 & 185.1 & 125.8 & 71.8 & 2.58$\times$ & 4/25/1 & 7/22/1 \\
\bottomrule
\end{tabular}
\end{table}

\subsection{Study 4: Component-wise Isolation of Fix~2 and Fix~3}
\label{sec:study_components}

Studies 1--3 compare FKS-CG-WS against KS-CG, which differs in two ways: it uses
flat-$k$ arc selection (Fix~2) and feasibility-preserving stage transitions
(Fix~3).
To isolate each fix, we compare against KS-CG-WS, which adds only the
feasibility-preserving transitions to KS-CG while retaining the $\gamma$-threshold.

\begin{itemize}
  \item KS-CG versus KS-CG-WS: measures Fix~3 (feasibility-preserving
        transitions) alone, with arc selection held constant.
  \item KS-CG-WS versus FKS-CG-WS: measures Fix~2 (flat-$k$ arc selection) alone,
        with feasibility-preserving transitions held constant.
\end{itemize}

Results for p1000-4000 Test Bed~A (30 instances) will be reported for this
ablation because its size and capacity regime produce clear quality differences
in Study~2.
The comparison uses the same instance set, CG LP output, and 300-second
restricted-MILP limit as the main study.
We keep this ablation separate from Tables~\ref{tab:kscg_fkscg}
and~\ref{tab:end_to_end} because it requires the additional KS-CG-WS control
run; those runs are not mixed into the aggregate tables until the complete
family is available.

\subsection{Study 5: Comparison Against Proven Optima}
\label{sec:study_proven_optima}

The final study evaluates solution quality against the best available results
from the literature for exact methods.
For Test Bed~A, we use the proven optimal values reported by
\citet{caserta2020corridor}, who solved 110 of the 150 instances to optimality.
For Test Bed~C, we use the proven optimal values and best-known upper bounds
from \citet{avella2021weakflow}, who provide the current state-of-the-art for
these harder instances.
The gap is computed as $\text{gap}(\%) = 100 \times (z^H - z^*) / z^*$, where
$z^*$ is the proven optimal or best-known objective.

Table~\ref{tab:proven_optima_gap} reports the average and maximum gaps for the
110 Test Bed~A instances with known optima.
The results show that FKS-CG-WS is exceptionally effective, finding the proven
optimal solution on 109 of 110 instances.
The single non-optimal solution (for instance p2000-2000-17) has a gap of just
0.0007\%, corresponding to an objective value less than one cost unit from the
optimum.
This confirms that the FKS design not only accelerates the search but also
consistently guides it to solutions of very high quality.

In contrast, the baseline KS-LP and KS-CG methods, while still strong, show
small but consistent gaps across all families.
Notably, KS-CG produces a large outlier gap of 0.27\% on one instance
(p1000-4000-17), where FKS-CG-WS finds the optimum. This suggests the flat-$k$
arc selection criterion is more robust than the $\gamma$-threshold.

\begin{table}[htbp]
\centering
\caption{Gap (\%) versus proven optimal solutions on 110 Test Bed~A instances.
         FKS-CG-WS finds the proven optimum on 109 of 110 instances.}
\label{tab:proven_optima_gap}
\footnotesize
\setlength{\tabcolsep}{4pt}
\begin{tabular}{lrrrrrr}
\toprule
& \multicolumn{2}{c}{KS-LP} & \multicolumn{2}{c}{KS-CG} & \multicolumn{2}{c}{FKS-CG-WS} \\
\cmidrule(lr){2-3}\cmidrule(lr){4-5}\cmidrule(lr){6-7}
Family (N) & Avg.\ gap & Max gap & Avg.\ gap & Max gap & Avg.\ gap & Max gap \\
\midrule
p800-4400 (21)   & 0.0021\% & 0.0246\% & 0.0010\% & 0.0055\% & 0.0000\% & 0.0000\% \\
p1000-1000 (30)  & 0.0018\% & 0.0373\% & 0.0001\% & 0.0025\% & 0.0000\% & 0.0003\% \\
p1000-4000 (21)  & 0.0029\% & 0.0173\% & 0.0211\% & 0.2726\% & 0.0000\% & 0.0000\% \\
p1200-3000 (19)  & 0.0027\% & 0.0268\% & 0.0014\% & 0.0096\% & 0.0000\% & 0.0000\% \\
p2000-2000 (19)  & 0.0026\% & 0.0136\% & 0.0012\% & 0.0088\% & 0.0000\% & 0.0007\% \\
\midrule
\textbf{Overall (110)} & \textbf{0.0024\%} & \textbf{0.0373\%} & \textbf{0.0047\%} & \textbf{0.2726\%} & \textbf{0.0000\%} & \textbf{0.0007\%} \\
\bottomrule
\end{tabular}
\end{table}

\subsection{Validating the KS-LP Baseline Against Published Results}
\label{sec:literature}

\paragraph{Why a re-implementation is necessary.}
The computational environments in G\&S~2012 and this study differ in solver
(CPLEX~12.2 versus Gurobi), hardware (Intel Xeon 2.27~GHz versus modern
hardware), and per-subproblem time limit (900~seconds versus 300~seconds).
These differences make direct comparisons of running times meaningless, and
they mean we cannot use the published G\&S timing or solution results as the
control in our ablation study.
Instead, we implement a faithful G\&S-style baseline --- same formulation, same
parameter choices ($\gamma$, $l_{\text{buck}}$, $\bar{N}_B$, $p$), same
kernel construction and bucket logic --- and run it as KS-LP in our own
environment.
All proposed variants (KS-CG, KS-CG-WS, FKS-CG-WS) are then compared against
this same KS-LP baseline, ensuring that solver, hardware, and time limit are
identical across every comparison in the study.

\paragraph{How we know the re-implementation is fair.}
The validity of this design depends entirely on whether our KS-LP recreation
is a faithful and adequate representation of the original algorithm.
If our re-implementation were systematically weaker than G\&S's, the study
would be unfairly easy: our new methods would beat a handicapped baseline.
We therefore validate KS-LP against G\&S's published results on the p1000-1000
Test Bed~A family, using the same benchmark quantity G\&S report: improvement
over the Avella et al.~\citep{avella2008effective} published upper bounds.

G\&S report an average improvement of $-0.13\%$ over those bounds, with 26
improvements out of 30 instances.
Our KS-LP achieves the same average improvement figure, improves 28 upper
bounds, and ties the remaining 2 --- despite using only one-third of
their per-subproblem time budget.
Because G\&S report rounded optimality gaps rather than raw integer objectives,
we also perform a secondary reconstruction check using
\[
  \tilde{z}^{H}_{\mathrm{BKS}}
  =
  \frac{z^{LB}_{\mathrm{BKS}}}
       {1-\mathrm{Gap}_{\mathrm{BKS}}/100},
\]
which gives the same aggregate conclusion: our KS-LP objective is slightly
lower than the reconstructed B-KS average (treated as an approximate
consistency check only, given the rounded inputs).

Our KS-LP re-implementation therefore produces \emph{at least as good}
solution quality as the published G\&S results on this family.
This confirms that KS-LP is a fair, non-handicapped representation of
the original algorithm, and that any gains attributed to KS-CG, KS-CG-WS,
or FKS-CG-WS in the following studies represent genuine improvements over
a credible baseline.
Detailed per-instance results appear in Appendix~\ref{app:detail}.
